\subsection{All unramified sections of the full, metabelian and single-cusp marked groups}

\begin{statement}{E.41}{Proved}{Complete compatible coordinates for all sections unramified outside three}
Use the actual based cubic action and the meaning of unramified
outside three in E.34. Let $Q=B_0$ be the rational cusp over one,
let $q$ be its positive peripheral generator, and let $I_Q$ be
its closed normal closure in $N$. All commutator subgroups below
are closed. Consider the three actual arithmetic quotients
\[
 W_{\rm free}=N,\qquad W_{\rm met}=N/N'',\qquad
 W_{\rm mark}=N/\gamma_3I_Q.
 \tag{E41.1}
\]
For each of these $W$, put
\[
 W_r=W/\gamma_{r+1}W,\qquad
 \mathcal A_r(W)=\gamma_rW/\gamma_{r+1}W.
 \tag{E41.2}
\]
These graded modules are finite free over $\mathbb Z_3$.
Let $T_r(W)$ and $J_r(W)$ be the matrices induced on them
by the two explicit matrices $A_\tau,C_V$ of E.37. Define
the following explicitly computable compact abelian group:
\[
 \mathcal C_r(W)=
 \operatorname{coker}\bigl((T_r(W)-1)(1+J_r(W)):
              \mathcal A_r(W)\longrightarrow\mathcal A_r(W)\bigr).
 \tag{E41.3}
\]
There are compatible noncanonical homeomorphisms of pointed
profinite spaces
\[
 \begin{split}
 \operatorname{Sect}_{\{3\}}(W_r)/W_r
       &\simeq\prod_{i=1}^r\mathcal C_i(W),\\
 \operatorname{Sect}_{\{3\}}(W)/W
       &\simeq\varprojlim_r\operatorname{Sect}_{\{3\}}(W_r)/W_r
        \simeq\prod_{i\ge1}\mathcal C_i(W).
 \end{split}
 \tag{E41.4}
\]
Here $\operatorname{Sect}_{\{3\}}$ consists of sections whose
geometric cocycle, relative to the specified tangent, factors
through $G_{\mathbb Q,\{3,\infty\}}$. The statement computes
this entire section family, including all its higher
coordinates. The product structures arise from specified
Hall lifts and Smith representatives; no group structure
on the section set is asserted.

The integral graded modules and their matrix bases are
\[
\begin{array}{c|c|l}
W&\mathcal A_r(W)&\text{basis for the finite calculation}\\ \hline
W_{\rm free}&\operatorname{Lie}_r(V)&\text{Hall commutators}\\
W_{\rm met}&(\operatorname{Lie}(V)/\operatorname{Lie}(V)'')_r
 & [x_i,x_j,x_{k_3},\ldots,x_{k_r}],\ i>j\le k_3\le\cdots\le k_r\\
W_{\rm mark}&
 (\operatorname{Lie}(V)/\gamma_3\mathfrak i_Q)_r
 &\text{Hall commutators with at most two occurrences of }q .
\end{array}
 \tag{E41.5}
\]
In the last row $\mathfrak i_Q$ is the Lie ideal generated by
the actual vector $e_Q$. Choose the other nine generators
to complete $q$ to a free pro-three basis inducing the
integral $V$-basis of E.37. The marking has the single
rational cusp $Q$ throughout this statement.

The first three factors in (E41.4) agree for all three
quotients and are, respectively,
\[
\begin{split}
 \mathcal C_1&=\mathbb Z_3^6\oplus(\mathbb Z/3)^3,\\
 \mathcal C_2&=\mathbb Z_3^{24}\oplus(\mathbb Z/3)^9
                                      \oplus(\mathbb Z/9)^5,\\
 \mathcal C_3&=\mathbb Z_3^{166}\oplus(\mathbb Z/3)^7
                  \oplus(\mathbb Z/9)^{56}
                  \oplus(\mathbb Z/27)^{25}
                  \oplus(\mathbb Z/81)^3.
\end{split}
 \tag{E41.6}
\]
For every degree the ranks and traces below, together
with the matrix (E41.3), evaluate all remaining factors.
Writing $n_r=\operatorname{rank}\mathcal A_r$ and
$t_r=\operatorname{tr}(J_r)$, the free rank is
$(n_r-t_r)/2$, and the torsion is killed by
\[
 3^{r+b_r},\qquad b_r=\max_{r\le j\le2r}(1+v_3(j)).
 \tag{E41.7}
\]
Thus a finite Smith calculation modulo $3^{r+b_r+1}$
determines all elementary divisors; there is a proved
bound on the precision required. In degree four the
actual matrices have the following sizes and free ranks:
\[
\begin{array}{c|r|r|r}
 W&n_4&t_4&\operatorname{rank}_{\mathbb Z_3}\mathcal C_4\\ \hline
 W_{\rm free}&2475&-21&1248\\
 W_{\rm mark}&2466&-22&1244\\
 W_{\rm met}&1485&-3&744
\end{array}
 \tag{E41.8}
\]
The supplied program constructs these integral matrices
and verifies the split-lattice certificates behind (E41.7).

\emph{Manuscript: \S51.16, folios 733v--736v, scans 172--178,
(371)--(379), Part E(c),(d). These are complete coordinates
for the family unramified outside three. The ramified
classes excluded in E.31 and E.38 remain outside this
family; their exact quotient maps are retained.}
\end{statement}
\begin{proof}
\emph{The based arithmetic quotient and an actual section parameter.}
The action on $N$ factors through
$\Gamma=P\rtimes\langle c\rangle$ by E.34, where
$P$ is free pro-three on $\tau,\tau'$ and
$c\tau c=\tau'$. E.37 proves that the two actual
conjugate Frobenii above seven form this free basis.
The same action descends to each group in (E41.1):
derived subgroups are characteristic, and the rational
cusp makes $I_Q$ stable under the based Galois action.

Every cocycle of $G_{\mathbb Q,\{3,\infty\}}$ with
values in any of these $W$ factors through $\Gamma$.
Indeed, over $K$, the map
$g\mapsto(a(g),\rho(g))$ has image in the pro-three
semidirect product $W\rtimes\rho(G_{K,\{3\}})$.
The universal property defining $P$ kills the same
kernel as in E.34. This supplies the actual comparison
from the arithmetic section problem to the following
explicit group problem.

Normalize a cocycle by geometric conjugation so that
$a(c)=1$. Coprime conjugacy supplies this normalization:
in each central graded layer the required correction
is obtained by division by two. The corrections lift
compatibly through the lower-central quotients, giving
the normalization in $W$ itself. The completeness and
torsion-freeness used here are proved below.
A normalized cocycle is freely specified by a single
element
\[
 x=a(\tau)\in W,\qquad a(\tau')=c(x),\qquad a(c)=1.
 \tag{E41.9}
\]
To verify existence, map the two free generators of
$P$ to $(x,\tau)$ and $(c(x),\tau')$ in
$W\rtimes P$. Freeness gives a continuous homomorphism
whose projection to $P$ is the identity. Conjugation
by $(1,c)$ interchanges these generator images, so the
homomorphism extends to a section over $\Gamma$.
Conversely the cocycle identities force these formulas.
The remaining geometric conjugations are exactly
\[
 W^{\langle c\rangle},\qquad
          x\longmapsto w x\tau(w)^{-1}.
 \tag{E41.10}
\]
In particular the normalized cocycle space is the
compact space $W$, and its section-class space is
the Hausdorff compact orbit space for this displayed
continuous compact-group action.

\emph{The integral metabelian comparison, including its topology.}
For the full free group, the integral graded commutator
isomorphism is the one in E.40 and
\cite[Theorem 2.1]{E_KochKukkukLabute}.
We prove the metabelian passage directly in the
pro-three setting. Set
\[
 R=\mathbb Z_3[[N]],\quad
 \Lambda=\mathbb Z_3[[N/N']]
      =\mathbb Z_3[[X_1,\ldots,X_{10}]],\quad
 \mathfrak m=(X_1,\ldots,X_{10}),
 \tag{E41.11}
\]
where $X_i=x_i-1$. The ring $R$ is the noncommutative
power-series ring in these ten variables with its
$(3,X_1,\ldots,X_{10})$ topology. Here is a direct
topological justification. The units $1+X_i$ in that
complete ring define a homomorphism from the free
pro-three group $N$, because their images at each
finite coefficient-and-degree level generate a
finite three-group. In the reverse direction,
$X_i\mapsto x_i-1$ defines a continuous power-series
map into $R$: in the group algebra of every finite
three-group over $\mathbb F_3$ the augmentation ideal
is nilpotent, and lifting to coefficients modulo
$3^a$ preserves nilpotence of the combined ideal.
The resulting maps are inverse on the generators
and therefore inverse on the complete rings.
The same argument after imposing commutativity
gives the topology on $\Lambda$ in (E41.11).

Unique decomposition of a nonempty word by its last
letter shows that the augmentation ideal of $R$ is
free as a left $R$-module on the $X_i$. Thus
\[
 0\longrightarrow R^{10}
     \xrightarrow{(a_i)\mapsto\sum a_iX_i}R
     \longrightarrow\mathbb Z_3\longrightarrow0
 \tag{E41.12}
\]
is an actual completed free resolution. Completed
tensoring with the right $R$-module $\Lambda$ gives
the exact identification
\[
 N'/N''\xrightarrow{\ d\ }\mathcal K
 :=\ker\bigl(\Lambda^{10}\xrightarrow{(X_i)}\Lambda\bigr).
 \tag{E41.13}
\]
For clarity, the continuous homology identification
in this assertion is the induced-module comparison
$\Lambda=\mathbb Z_3\widehat\otimes_{\mathbb Z_3[[N']]}R$.
There is a continuous coset section
$(a_1,\ldots,a_{10})\mapsto x_1^{a_1}\cdots x_{10}^{a_{10}}$;
the usual bar contraction along this transversal
identifies the first homology after tensoring
(E41.12) with $H_1(N',\mathbb Z_3)=N'/N''$.
On a loop its chain map is precisely the continuous
Fox derivative
\[
 d(x_i)=e_i,\qquad d(uv)=d(u)+u^{\rm ab}d(v),\qquad
       \sum X_i d_i(u)=u^{\rm ab}-1.
 \tag{E41.14}
\]
This describes the isomorphism, rather than only
its two objects. Both sides of (E41.13) are compact
and the target is a closed kernel. Its continuous
bijection is consequently a homeomorphism. In
particular the topology on $N'/N''$ is exactly the
induced $(3,\mathfrak m)$ topology in (E41.13).

The elements $X_i$ form a regular sequence in
$\Lambda$. Its Koszul complex presents $\mathcal K$
by the vectors $X_ie_j-X_je_i$ and their Jacobi
relations. Commutation with $x_i$ in the group
corresponds under (E41.14) to multiplication by
$-X_i$. It follows, with closed images throughout,
that
\[
 \gamma_r(W_{\rm met})=\mathfrak m^{r-2}\mathcal K
 \quad(r\ge2),\qquad
 \mathcal A_r(W_{\rm met})
  =\ker\bigl(\operatorname{Sym}^{r-1}V\otimes V
                        \longrightarrow\operatorname{Sym}^rV\bigr).
 \tag{E41.15}
\]
The filtration comparison used in the second
equality is, explicitly,
\[
 \mathcal K\cap\mathfrak m^{k+1}\Lambda^{10}
                   =\mathfrak m^k\mathcal K\qquad(k\ge0).
\]
One inclusion follows from the degree-one
generators. For the other, decompose a syzygy
into its homogeneous parts of degrees at least
$k+1$. Homogeneous Koszul exactness expresses
each part by $X_ie_j-X_je_i$ with coefficients
of degree at least $k$. Assemble these
coefficients into power series; their convergent
sum is the required element of
$\mathfrak m^k\mathcal K$. Thus the second
equality of (E41.15) is exactly the homogeneous
degree-$r-1$ syzygy map. This also proves that
no extra initial relations or completion
torsion occur.
This is the integral free metabelian Lie quotient;
its commutator-to-Koszul map is the one in
\cite[Theorem 6.2, including (6.3)--(6.4)]{E_PapadimaSuciuChen}.
Its stated Chen basis is the integral basis in
\cite[\S2.4, equation(2), p.~5]{E_SatohChen}.
The preceding resolution proves explicitly the
passage to the pro-three completion used here.

\emph{The exact single-cusp marked comparison.}
Write the other nine degree-one generators as $X$.
Integral Lazard elimination gives
\[
 \operatorname{Lie}(X\oplus\mathbb Z_3q)
       =\operatorname{Lie}(X)\ltimes
          \operatorname{Lie}\bigl(\mathbb Z_3q\otimes T(X)\bigr),
 \tag{E41.16}
\]
where the free generators of the ideal are
$[q,x_{i_1},\ldots,x_{i_s}]$, $s\ge0$.
The version over an arbitrary commutative ring is
opened in \cite[\S2, Lazard Elimination Theorem,
p.~3]{E_StohrLazard}. Each ideal generator contains
one $q$. Its third lower-central term is therefore
exactly the direct sum of the multihomogeneous
components containing at least three $q$'s.

We verify that this is the initial ideal of the
actual group subgroup $\gamma_3I_Q$. In the Magnus
power-series expansion, $u-1$ for $u\in I_Q$ has
at least one occurrence of the variable $q-1$
in every monomial: $I_Q$ is the kernel of the
retraction setting $q=1$. The expansion of an
element of $\gamma_3I_Q$ consequently has at least
three occurrences in each monomial. Taking the
first nonzero ordinary degree gives one inclusion
in the proposed initial ideal. In the reverse
direction each generator of the ideal in
(E41.16) is the leading term of the corresponding
group commutator in $I_Q$. Every bracket of at
least three such generators is the leading term
of an element of $\gamma_3I_Q$. Products and
$\mathbb Z_3$ powers realize their sums and
scalar multiples. This proves the reverse
inclusion in every degree, including the closure.
Hence the graded quotient is exactly the last
row of (E41.5). The Hall basis is multihomogeneous;
discarding its elements with at least three $q$'s
leaves a free integral basis. In particular the
marked quotient has no extra graded torsion.

All three descriptions are functorial for the
relevant based arithmetic action. In the marked
case the line $\mathbb Z_3e_Q$ is stable; the chosen
complement need not be stable, and the quotient
matrix is obtained by actual substitution followed
by the just-proved ideal reduction. Thus the
displayed matrices use the whole rank-ten action.

\emph{Every layer, its free rank and a finite torsion computation.}
The graded modules in (E41.5) are quotients of
degree-$r$ Lie powers with the actual induced
arithmetic action. Their rational Frobenius
eigenvalues have absolute values $p^{j/2}$ with
$r\le j\le2r$. In particular none is one.
An element fixed by a section-twisted action on
$W_r$ would have a first nonzero graded coordinate
fixed by the original action: inner conjugation
acts trivially on the associated graded module.
The preceding eigenvalue calculation excludes
such a coordinate. We have therefore proved
\[
 H^0(\Gamma,W_r^{\mathrm{section}})=1
              \quad\text{for every section and every }r.
 \tag{E41.17}
\]

The central action on $\mathcal A_r$ is unaffected
by a lower section. A normalized additive cocycle
is freely specified by its value at $\tau$, and
the normalized coboundaries are exactly
$(T_r-1)\mathcal A_r^{\langle c\rangle}$.
The image of $1+J_r$ equals that fixed lattice
because two is a unit. This proves (E41.3) as an
evaluation map and shows that it is the complete
parameter group for each nonempty central fiber.
Frobenius at seven makes $T_r-1$ rationally
injective, so the free rank is the minus-space
dimension $(n_r-t_r)/2$.

Here are formulas evaluating those dimensions in
every degree. For the full group they are (E40.5).
For the metabelian group and $r\ge2$, put
$h_s=[t^s](1-t)^{-4}(1+t)^{-6}$. Taking the trace
in the exact kernel (E41.15) gives
\[
 n_r^{\rm met}=(r-1){r+8\choose r},\qquad
 t_r^{\rm met}=-2h_{r-1}-h_r.
 \tag{E41.18}
\]
For the marked group define
\[
 w_r(9)=\frac1r\sum_{d\mid r}\mu(d)9^{r/d},\qquad
 u_r=\frac1r\sum_{d\mid r}\mu(d)
                  (4+(-1)^d5)^{r/d}.
\]
Complex conjugation on $V/\mathbb Z_3e_Q$
has plus and minus multiplicities four and five;
on $q$ it is minus one. Averaging a projection
by $c$ supplies an invariant integral complement
for computing this trace. Elimination separates
zero, one and two $q$ occurrences. Their dimensions
and traces give, for $r\ge2$,
\[
\begin{split}
 n_r^{\rm mark}
  &=w_r(9)+9^{r-1}
       +\frac{(r-1)9^{r-2}
                       -\mathbf1_{2\mid r}9^{r/2-1}}2,\\
 t_r^{\rm mark}
  &=u_r+(-1)^r
       +\frac{(r-1)(-1)^r
                       -\mathbf1_{2\mid r}9^{r/2-1}}2.
\end{split}
 \tag{E41.19}
\]
For example the last terms are obtained from
$\bigwedge^2(\mathbb Z_3q\otimes T(X))$; its
degree-$r$ trace is one half of the trace on the
tensor square minus the diagonal trace of $c^2$.
Witt's formula handles the zero-$q$ part. In
degree one each group has rank ten and trace
minus two. These formulas prove (E41.8).

The bound (E41.7) passes through these exact
integral quotients. Namely, the actual split
sublattice $V_0=T_3T_D+j(T_3E)$ contains $3V$.
The image of the corresponding degree-$r$ lattice
contains $3^r\mathcal A_r(W)$. It is a direct sum
by the number of elliptic letters: the metabelian
relations are homogeneous for this count, and
the marked letter is toric. Frobenius at the
$K$-prime above two acts on each summand by
$(-2)^j$ with $r\le j\le2r$. The identity
$v_3((-2)^j-1)=1+v_3(j)$ proves
\[
 3^{r+b_r}(F-1)^{-1}\mathcal A_r(W)
                         \subset\mathcal A_r(W).
 \tag{E41.20}
\]
Since the rational invariant space is zero,
the torsion of the cokernel (E41.3) is precisely
$H^0(\Gamma,\mathcal A_r\otimes\mathbb Q_3/\mathcal A_r)$.
Equation (E41.20) kills this group by the stated
power of three. Substituting the two known
ten-dimensional matrices into the explicit
bases (E41.5), reducing by antisymmetry, Jacobi
and the displayed ideal, and taking Smith form
at the stated precision therefore evaluates
every factor by a finite procedure. Exact Smith
representatives may be chosen over
$\mathbb Z_{(3)}$ since the input matrices are
integral; unit denominators define continuous
maps on the resulting $\mathbb Z_3$ coordinates.

The first three evaluations are E.34, the
unramified summand of E.37, and E.40. The two
defining group kernels in (E41.1) lie in
$\gamma_4N$, as proved in (E40.13)--(E40.14),
so their actual graded comparisons identify
these three layers. This proves (E41.6).
The program \texttt{e\_unramified\_sections.py}
also performs the first distinct degree, four.
It constructs exact Hall and Chen action
matrices, the marked matrix by actual quotient
reduction, and the matrices of the maps
$B,C$ with $BC=3I$ from E.40.
Its saved integer identities verify
$B_rC_r=3^rI$, conjugation squared equal to the
identity, the traces in (E41.8), and Frobenius
conjugacy to the split letter-count blocks.
Every operation uses exact Python integers.
The output includes the sparse matrices and
rank tables through degree ten.

\emph{Compatible normal forms and the inverse limit.}
We now construct (E41.4); central liftability
alone would not supply these coordinates.
Choose the ordered commutator representatives
for each basis in (E41.5). Their $\mathbb Z_3$
powers, ordered by degree and then by the
chosen basis order, give continuous unique
coordinates on each $W_r$. This follows
inductively from its central kernel
$\mathcal A_r$, and supplies a continuous set
section
\[
 h_r:W_{r-1}\longrightarrow W_r
\]
by setting the last-degree coordinates to zero.
Choose Smith-coordinate representatives
$s_r:\mathcal C_r(W)\to\mathcal A_r(W)$,
continuous and taking zero to zero: retain
$\mathbb Z_3$ in a free coordinate and the
digits $0,\ldots,3^e-1$ in a factor
$\mathbb Z/3^e$. Starting with $x_0=1$, define
\[
 x_r=h_r(x_{r-1})\,s_r(z_r),
       \qquad z_r\in\mathcal C_r(W).
 \tag{E41.21}
\]
The additive element $s_r(z_r)$ is viewed here
in the central kernel. Formula (E41.9) makes
each $x_r$ an actual normalized section.
In particular every lower section class has
a lift in the unramified family.

We prove uniqueness of these representatives
inductively. A normalized section at level
$r$ has a lower representative at level
$r-1$ already classified by the induction.
There is a unique element of
$W_{r-1}^{\langle c\rangle}$ aligning it
with that representative: two aligning
elements differ by the stabilizer (E41.17).
This element has a lift in
$W_r^{\langle c\rangle}$. To see the latter
surjectivity explicitly, lift it arbitrarily
to $\widetilde w$ and write
$\delta=c(\widetilde w)\widetilde w^{-1}$
additively in the central kernel. Then
$c\delta=-\delta$, and
$(\delta/2)\widetilde w$ is fixed by $c$.
After this alignment the
remaining difference lies in $\mathcal A_r$.
Its changes under the remaining central
fixed conjugators are exactly
$(T_r-1)\mathcal A_r^{\langle c\rangle}$.
The chosen Smith representative therefore
gives one and only one $z_r$. Equation
(E41.17) again shows uniqueness of the full
aligning conjugator. This proves the finite
product bijection, with the asserted
compatibility under projection.

For completeness, every pro-three group here
is its inverse limit of lower-central
quotients. An element in every
$\gamma_rW$ dies in every finite
three-group quotient, since such a quotient
is nilpotent; it is therefore the identity.
Compactness proves surjectivity onto the
inverse limit by the nested cosets of the
closed subgroups $\gamma_rW$. Hence
\[
 W\xrightarrow{\sim}\varprojlim_r W_r.
 \tag{E41.22}
\]
The words $x_r$ in (E41.21) have a compatible
limit $x\in W$, and (E41.9) gives an actual
continuous full section. For any other
normalized section the unique finite-level
aligning conjugators are compatible; their
limit lies in $W^{\langle c\rangle}$.
This proves the full bijection in (E41.4),
including injectivity on geometric-conjugacy
classes. The construction is continuous
because every coordinate at a finite stage
uses continuous Hall lifts and Smith
representatives. A continuous bijection
from the compact product to the Hausdorff
orbit space (E41.10) is a homeomorphism.

Finally the comparisons among the three
families are the induced maps from the
actual epimorphisms
\[
 N\longrightarrow N/N'',\quad\ker=N'',
 \qquad
 N\longrightarrow N/\gamma_3I_Q,\quad\ker=\gamma_3I_Q.
 \tag{E41.23}
\]
At level $r$ their kernels are respectively
$N''\gamma_{r+1}N/\gamma_{r+1}N$ and
$\gamma_3I_Q\,\gamma_{r+1}N/\gamma_{r+1}N$.
On the section coordinates, apply the
displayed quotient to the word $x_r$ and
then the just-proved unique normalization
in its target. These formulas specify the
comparison even when the independently
chosen Smith coordinates do not make it
a coordinate projection. They also retain
the common quotient $N/(N''\gamma_3I_Q)$
and its two quotient maps. The full
ramified section problem is not replaced
by this unramified family: E.39 gives its
additional inertia coordinates and their
actual correction words, while E.38
exhibits a ramified central obstruction.
\end{proof}
